% Evaluation opening: replace the build-description paragraph.
\noindent\textbf{Builds and measurement.}
The current-build series is based on \texttt{f856c7b}: three
100-hop runs per continuation mode and three runs of the
representative paused-adaptation treatment P. All nine runs
share the same node build. Additional timing probes are confined
to \texttt{payctl} and are enabled in all six chain runs, without
changes to production node paths. The complete N/R/C/P matrix
on \texttt{721800c} remains archived; historical-reader measurements
retain \texttt{721800c}, and the injected-latency study retains
\texttt{9879f64}. Chain summaries use medians and ranges of three
run totals; mixed-load summaries use medians of three run-level
statistics.

% Replace the continuous-spending subsection text.
\subsection{Continuous Spending Without Per-Hop Settlement}

Three independent NoSync wallets rotate 100~CAL through online
100-hop chains, using synchronous member commits and user-funded
final FUEL. Fast mode proceeds after certified receipt; wait mode
also requires the driver's observation of wallet-verified public
execution between hops. Both measure
$\mathsf{Sent}_1 \rightarrow \mathsf{Ready}_{100}$, including
successor construction and retries. Wait mode adds 99 intermediate
waits, not the final payment's public confirmation.

Median chain times are 4.707\,s and 64.479\,s, respectively,
a 13.70-fold ratio of medians; run ranges are 4.700--4.726\,s
and 64.387--64.670\,s. Fast uses 300 attempts for 300 payments,
whereas wait uses 414. All 297 fast successors use the preceding
TXCer and are first sent before the earliest parent
application-commit completion across the four replicas, with a
minimum lead of 183.439\,ms. Actual child-before-source obligations
number 2, 4, and 1 in the three fast runs; early sending alone is
not counted as out-of-order execution. All 600 payments close and
pass the accounting, input, replica, fee, and outbox audits.
Fast-mode all-payment member closure has a median of 5.526\,s
from the first send. Each chain spends 9,400~FUEL. For the fixed
transaction shape, the TXCer is 779~B and requests are 1,979~B
initially and 2,766~B thereafter.

\noindent\textbf{Client timing.}
Records are matched by receiving wallet, block height, and payment
fact. This decomposition uses application-commit completion at
\texttt{committee0}, the fetch source, rather than the earliest
replica used in the lead check above. For height $h$,
\texttt{BlockData} first obtains the signed header at $h+1$ and
then fetches the block and results at $h$. Commit-to-complete-fetch
therefore remains a joint interval covering successor evidence,
replica progress, polling, and retrieval, not pure query overhead.
Its wait-mode P50/P95 are 285.884/314.509\,ms. The separately
observed post-fetch verification, preparation/local-recording,
and post-transaction-callback-to-driver intervals have P95s of
0.236, 1.343, and 4.713\,ms, respectively. The verification
checkpoint is the \texttt{PrepareBlock} entry after successful
\texttt{VerifyBlock}; the callback follows local transaction
completion, not a durability flush.

Cross-process intervals use same-host Unix-nanosecond timestamps;
whole-chain durations use a monotonic clock. Driver polling,
follower polling, and the commit setting remain 5, 25, and
250\,ms. Fast successors can precede public observation, so
negative observation-to-next-build differences record overlap,
not negative waiting time. The chain ratio compares the measured
end-to-end modes, including retries and evidence observation.

% Replace the mixed-load subsection text.
\subsection{Source Resolution Under Representative Mixed Load}

The fourteen-service P treatment schedules 7,000 ordinary payments
at 100 payments/s over 70\,s, alongside eight independent
cross-organization parent--child pairs, using authorized
organization-funded FUEL. Parent publication and \texttt{INSTALL}
are withheld, automatic source resubmission is disabled, and all
adaptation endpoints are paused; consensus and compensation
decisions remain active. Each late source is released after its
compensation is observed. Ordinary traffic spans the injected
fault and recovery intervals.

All 21,048 payments across the three runs close, including
21,000 ordinary payments and 24 parent--child pairs. Each run
pays and recovers 800~CAL. Before adaptation resumes, all four
committee replicas report all eight targets \texttt{Recovered},
with their representation \texttt{Committed} and
\texttt{Materialized} flags both false. These 96 replica--target
checks establish compensation and repayment before adaptation
resumption in the measured executions. Within each run,
committee state hashes agree at a common height, and CAL/FUEL,
input, fee, reservation, and outbox audits pass. Per-run fees are
659,544~FUEL: 589,384 in rewards and 70,160 burned.

Table~\ref{tab:current-p} reports ordinary receipt latency together
with pre-send backpressure and drain. Median run-level receipt
P50/P95 are 54.538/147.549\,ms; public and member-observation P95s
are 2.187 and 2.255\,s. All three runs reach the 256-outstanding
limit. Dispatch P95 reaches 780.727\,ms and scheduled-to-receipt
P95 reaches 851.742\,ms; the latter is computed per payment,
not by adding quantiles. Actual send spans are 70.080--70.771\,s,
followed by 0.955--1.115\,s of ordinary-payment drain. This
representative P series demonstrates economic closure during
ordinary traffic with adaptation unavailable; it is not a
current-build comparison of interference across N/R/C/P.